\clearpage
\begin{explanationbox}
\textbf{\textcolor{CExplanation}{先认清：焦半径就是两条连向焦点的线段}}

取椭圆 $x^2/25+y^2/16=1$，$a=5,c=3,e=3/5$。图中 $P=(3,16/5)$ 确实在椭圆上。
\begin{center}
\begin{tikzpicture}[x=1cm,y=1cm,line cap=round,line join=round,font=\small]
\draw[black!25,rounded corners=3pt] (0,0) rectangle (14.4,9.4);
\node[font=\small\bfseries] at (7.2,8.9) {两条焦半径：蓝色连左焦点，橙色连右焦点};
\begin{scope}[shift={(4.7,4.7)},x=0.7cm,y=0.7cm]
\draw[black!35,->] (-5.7,0)--(5.3,0) node[right] {$x$};
\draw[black!35,->] (0,-4.5)--(0,4.6) node[above] {$y$};
\draw[black!65,thick] (0,0) ellipse (5 and 4);
\coordinate (P) at (3,3.2);
\coordinate (Fone) at (-3,0);
\coordinate (Ftwo) at (3,0);
\draw[blue!70!black,line width=1.1pt] (Fone)--(P);
\draw[orange!85!black,line width=1.1pt] (Ftwo)--(P);
\foreach \q in {P,Fone,Ftwo} \fill (\q) circle (1.6pt);
\node[above right] at (P) {$P(3,\frac{16}{5})$};
\node[below] at (Fone) {$F_1(-3,0)$};
\node[below] at (Ftwo) {$F_2(3,0)$};
\node[below left] at (0,0) {$O$};
\path (Fone)--(P) node[pos=.28,blue!70!black,above,sloped] {$r_1=\frac{34}{5}$};
\node[orange!85!black,left] at (3,1.1) {$r_2=\frac{16}{5}$};
\end{scope}
\node[anchor=north west,align=left] at (9.1,7.9) {$\dfrac{x^2}{25}+\dfrac{y^2}{16}=1$\\[1em]$a=5,\ c=3,\ e=\dfrac35$\\[1em]\textcolor{blue!70!black}{$r_1=a+ex_0$}\\[.5em]\textcolor{orange!85!black}{$r_2=a-ex_0$}\\[1em]$r_1+r_2=10=2a$};
\node at (7.2,.65) {“左加右减”指所连焦点的位置；$x_0$ 保留正负号，不换成 $|x_0|$。};
\end{tikzpicture}
\end{center}

\textbf{读图：} 蓝线连左焦点，长 $5+\frac35\cdot3=\frac{34}{5}$；
橙线连右焦点，长 $5-\frac35\cdot3=\frac{16}{5}$。
和为 $10$，与椭圆定义相符。

\textbf{推广到任意椭圆上的点：}
\[
\textcolor{blue!70!black}{r_1=a+ex_0},\qquad
\textcolor{orange!85!black}{r_2=a-ex_0},\qquad r_1+r_2=2a.
\]
这里的“加、减”对应\emph{焦点位置}；不是把两条线段的长短永远固定。
例如 $x_0<0$ 时 $ex_0<0$，左焦半径反而较短。
\medskip
\textbf{\textcolor{CExplanation}{补充直觉：两段长度和已知，为什么只设一个 $t$？}}
既然 $r_1+r_2=2a$，就把 $a$ 看作每条的“平均长度”：
\[
 \boxed{r_1=a+t,\qquad r_2=a-t.}
\]
这相当于从一条挪出 $t$ 给另一条，\emph{总和自动保持 $2a$}。
与分别设 $r_1,r_2$ 再求和、差相比，这里\textbf{只需确定 $t$ 一个未知量}；
而且 $t$ 可正可负，不限定 $P$ 在左侧还是右侧。
\end{explanationbox}
\clearpage
\begin{explanationbox}
\textbf{\textcolor{CExplanation}{再看变化：一条增、一条减，和保持不变}}
\begin{center}
\begin{tikzpicture}[x=1cm,y=1cm,line cap=round,line join=round,font=\small]
\draw[black!25,rounded corners=3pt] (0,0) rectangle (14.4,8.8);
\node[font=\small\bfseries] at (7.2,8.3) {从左端到右端：蓝色变长、橙色变短，总长度不变};
\begin{scope}[shift={(2.4,5.1)},x=.36cm,y=.36cm]
\draw[black!25] (-5.5,0)--(5.5,0);
\draw[black!60] (0,0) ellipse (5 and 4);
\coordinate (P) at (-5,0);
\coordinate (Fone) at (-3,0);
\coordinate (Ftwo) at (3,0);
\draw[blue!70!black,line width=2.3pt] (Fone)--(P);
\draw[orange!85!black,densely dashed,line width=.9pt] (Ftwo)--(P);
\foreach \q in {P,Fone,Ftwo} \fill (\q) circle (1.6pt);
\node[below] at (Fone) {$F_1$};
\node[below] at (Ftwo) {$F_2$};
\node[above] at (P) {$P$};
\end{scope}
\node at (2.4,3.2) {$x_0=-5$};
\node[blue!70!black] at (2.4,2.6) {$r_1=2$};
\node[orange!85!black] at (2.4,2.0) {$r_2=8$};
\fill[blue!45] (0.5999999999999999,1.35) rectangle (1.3199999999999998,1.7);
\fill[orange!50] (1.3199999999999998,1.35) rectangle (4.2,1.7);
\node at (2.4,.7) {$r_1r_2=16$};
\begin{scope}[shift={(7.199999999999999,5.1)},x=.36cm,y=.36cm]
\draw[black!25] (-5.5,0)--(5.5,0);
\draw[black!60] (0,0) ellipse (5 and 4);
\coordinate (P) at (0,4);
\coordinate (Fone) at (-3,0);
\coordinate (Ftwo) at (3,0);
\draw[blue!70!black,line width=2.3pt] (Fone)--(P);
\draw[orange!85!black,densely dashed,line width=.9pt] (Ftwo)--(P);
\foreach \q in {P,Fone,Ftwo} \fill (\q) circle (1.6pt);
\node[below] at (Fone) {$F_1$};
\node[below] at (Ftwo) {$F_2$};
\node[above] at (P) {$P$};
\end{scope}
\node at (7.199999999999999,3.2) {$x_0=0$};
\node[blue!70!black] at (7.199999999999999,2.6) {$r_1=5$};
\node[orange!85!black] at (7.199999999999999,2.0) {$r_2=5$};
\fill[blue!45] (5.3999999999999995,1.35) rectangle (7.199999999999999,1.7);
\fill[orange!50] (7.199999999999999,1.35) rectangle (9.0,1.7);
\node at (7.199999999999999,.7) {$r_1r_2=25$};
\begin{scope}[shift={(12.0,5.1)},x=.36cm,y=.36cm]
\draw[black!25] (-5.5,0)--(5.5,0);
\draw[black!60] (0,0) ellipse (5 and 4);
\coordinate (P) at (5,0);
\coordinate (Fone) at (-3,0);
\coordinate (Ftwo) at (3,0);
\draw[blue!70!black,line width=2.3pt] (Fone)--(P);
\draw[orange!85!black,densely dashed,line width=.9pt] (Ftwo)--(P);
\foreach \q in {P,Fone,Ftwo} \fill (\q) circle (1.6pt);
\node[below] at (Fone) {$F_1$};
\node[below] at (Ftwo) {$F_2$};
\node[above] at (P) {$P$};
\end{scope}
\node at (12.0,3.2) {$x_0=5$};
\node[blue!70!black] at (12.0,2.6) {$r_1=8$};
\node[orange!85!black] at (12.0,2.0) {$r_2=2$};
\fill[blue!45] (10.2,1.35) rectangle (13.079999999999998,1.7);
\fill[orange!50] (13.079999999999998,1.35) rectangle (13.8,1.7);
\node at (12.0,.7) {$r_1r_2=16$};
\node at (7.2,7.6) {同一椭圆 $x^2/25+y^2/16=1$，三图与下方长度条均保持同尺度。};
\end{tikzpicture}
\end{center}

端点处两条线段共线，用蓝色粗实线、橙色细虚线区分；长度条显示各自全长。
三幅真实状态直接检验“左加右减”。随着 $x_0$ 从 $-a$ 增到 $a$，
$r_1$ 从 $a-c$ 增到 $a+c$，$r_2$ 反向变化，每条都为正。

\textbf{长度条还告诉我们：} 两条之和相同，长度相等时乘积最大；在两端时乘积最小。
一般地
\[
r_1r_2=a^2-e^2x_0^2\in[b^2,a^2].
\]
最大值在 $x_0=0$（两个短轴端点）取得；最小值在 $x_0=\pm a$（两个长轴端点）取得。

\textbf{理解不等于证明：} 三态图用来理解、验算符号；下面用准线或勾股平方差证明\emph{所有点}都成立。
\end{explanationbox}
